\section{Introduction}

\begin{frame}{Outline}
  \begin{itemize}
    \item Theory
      \begin{itemize}
        \item What is the centrifugal wall?
        \item The ``spinnability'' criterion.
        \item The Helium solvation shell model and the renormalized rotational eigenstates.
      \end{itemize}
    \item Experiments
      \begin{itemize}
        \item We have spun several molecules in droplets: CS2, OCS, DIB, NO-dimer. 
        \item Compare the CS2 and OCS data from the centrifuge with higher and lower acceleration. 
        \item Results with a decelerating centrifuge. 
        \item Why a truncated cfg would be helpful.
      \end{itemize}
    \item Comparison to a new thermalization model from Vienna collaboration
    \item Next steps
      \begin{itemize}
        \item Decelerating cfg starts spinning cs2 at ~ 50 GHz. Let's shift the cfg central frequency higher --- how much faster can we spin CS2?
        \item The new cfg shaper!
      \end{itemize}
  \end{itemize}
\end{frame}
